\documentclass[10pt]{extarticle}
\usepackage[a4paper,margin=1.5cm]{geometry}
\usepackage[T1]{fontenc}
\usepackage{amsmath,amssymb}
\usepackage{graphicx}
\usepackage[export]{adjustbox}
\usepackage{multicol}
\usepackage{enumitem}
\usepackage{booktabs,array,tabularx}
\setlength{\columnsep}{0.7cm}
\setlength{\parindent}{0pt}
\setlength{\parskip}{3pt}
\pagestyle{empty}
\makeatletter
% figure registry (figdefs.tex is generated by the extract step)
\newcommand{\qfigdef}[4]{\@namedef{qf@file@#1}{#2}\@namedef{qf@w@#1}{#3}\@namedef{qf@wide@#1}{#4}}
\newcommand{\qf@check}[1]{\@ifundefined{qf@file@#1}{\PackageError{score2}{Unknown figure id #1}{}}{}}
\newcommand{\qf@img}[4][]{\includegraphics[width=#2\dimexpr\@nameuse{qf@w@#4}pt\relax,max width=#3,max height=0.45\textheight,#1]{\@nameuse{qf@file@#4}}}
% Full-width questions: finalize wraps every question that contains a wide figure in \qfullbegin...\qfullend,
% so the WHOLE question (stem, figures, options) is set across the page in reading order, then the two columns
% resume. The question's head - its number, stem and first figure (\qfullsplit marks the end of the head) -
% is kept on one page, so a number is never left at a page foot with its figure overleaf; the rest may break
% across pages normally, so no big gaps are left. Without a \qfullsplit, the whole question is kept together
% if it is under 45% of a page. Never write these yourself.
\newif\ifqf@full
\newif\ifqf@boxed
\newsavebox{\qf@box}
\newcommand{\qf@open}{\global\qf@fulltrue\par\global\qf@boxedtrue
  \setbox\qf@box\vbox\bgroup\hsize\textwidth\linewidth\textwidth}
\newcommand{\qf@emit}[1]{\par\egroup\global\qf@boxedfalse
  \ifdim\dimexpr\ht\qf@box+\dp\qf@box\relax<#1\textheight\noindent\usebox{\qf@box}\par\else\unvbox\qf@box\fi}
\newcommand{\qfullsplit}{\ifqf@boxed\qf@emit{0.85}\fi}
\newcommand{\qf@close}{\ifqf@boxed\qf@emit{0.45}\fi\global\qf@fullfalse}
\newcommand{\qfullbegin}{\end{multicols}\qf@open}
\newcommand{\qfullend}{\qf@close\begin{multicols}{2}}
\newcommand{\qfullstart}{\qf@open}
\newcommand{\qfulldone}{\qf@close}
% \qfig[scale]{id}: block figure, centred, directly where it is written (size set by the tool).
% Fallback if a wide figure is not inside a full-width question: break out of the columns in place - no floats.
\newcommand{\qfig}[2][1]{\qf@check{#2}%
  \ifqf@full\qf@blk{#1}{#2}\else
  \ifnum\@nameuse{qf@wide@#2}=1 \expandafter\@firstoftwo\else\expandafter\@secondoftwo\fi
  {\end{multicols}\begin{center}\qf@img{#1}{\textwidth}{#2}\end{center}\begin{multicols}{2}}%
  {\qf@blk{#1}{#2}}\fi}
\newcommand{\qf@blk}[2]{\par\smallskip{\centering\qf@img{#1}{\linewidth}{#2}\par}\smallskip}
% \qfigi[scale]{id}: inline figure (inside option lists / tables / minipages); never breaks columns.
\newcommand{\qfigi}[2][1]{\qf@check{#2}\qf@img[valign=c]{#1}{\linewidth}{#2}}
\makeatother
% structure macros - lint reads these to check completeness and figure placement
\newcommand{\papersec}[1]{\par\bigskip{\centering\Large\bfseries PAPER-#1\par}\medskip}
\newcommand{\partsec}[2]{\par\medskip\noindent{\large\bfseries PART-#1: #2}\par\smallskip}
\newcommand{\sect}[1]{\par\smallskip\noindent\textbf{#1}\par}
\newcommand{\qn}[1]{\par\smallskip\noindent\textbf{#1.}~}
\newcommand{\paraq}[2]{\par\smallskip\noindent\textbf{Paragraph for Questions #1 and #2}\par}
\newenvironment{qoptions}{\begin{itemize}[leftmargin=1.8em,itemsep=1pt,topsep=2pt,label={}]}{\end{itemize}}
\newcommand{\opt}[2]{\item[(#1)] #2}
% figure registry (generated by score2_tools.py finalize)
\qfigdef{p1_phy_s1_q01_1}{Score Advanced - Test 21 (08 Apr 2024)-img/p1_phy_s1_q01_1.png}{155.5}{0}
\qfigdef{p1_phy_s1_q02_1}{Score Advanced - Test 21 (08 Apr 2024)-img/p1_phy_s1_q02_1.png}{48.7}{0}
\qfigdef{p1_phy_s1_q03_1}{Score Advanced - Test 21 (08 Apr 2024)-img/p1_phy_s1_q03_1.png}{222.1}{0}
\qfigdef{p1_phy_s1_q04_1}{Score Advanced - Test 21 (08 Apr 2024)-img/p1_phy_s1_q04_1.png}{101.1}{0}
\qfigdef{p1_phy_s1_q05_1}{Score Advanced - Test 21 (08 Apr 2024)-img/p1_phy_s1_q05_1.png}{106.3}{0}
\qfigdef{p1_phy_s1_q06_1}{Score Advanced - Test 21 (08 Apr 2024)-img/p1_phy_s1_q06_1.png}{161.3}{0}
\qfigdef{p1_phy_s2_para7_8_1}{Score Advanced - Test 21 (08 Apr 2024)-img/p1_phy_s2_para7_8_1.png}{139.8}{0}
\qfigdef{p1_phy_s2_para9_10_1}{Score Advanced - Test 21 (08 Apr 2024)-img/p1_phy_s2_para9_10_1.png}{182.2}{0}
\qfigdef{p1_phy_s3_para1_2_1}{Score Advanced - Test 21 (08 Apr 2024)-img/p1_phy_s3_para1_2_1.png}{127.1}{0}
\qfigdef{p1_phy_s3_para3_4_1}{Score Advanced - Test 21 (08 Apr 2024)-img/p1_phy_s3_para3_4_1.png}{106.0}{0}
\qfigdef{p1_phy_s4_q08_1}{Score Advanced - Test 21 (08 Apr 2024)-img/p1_phy_s4_q08_1.png}{182.2}{0}
\qfigdef{p1_phy_s4_q09_1}{Score Advanced - Test 21 (08 Apr 2024)-img/p1_phy_s4_q09_1.png}{136.1}{0}
\qfigdef{p1_chem_s1_q03_A}{Score Advanced - Test 21 (08 Apr 2024)-img/p1_chem_s1_q03_A.png}{51.4}{0}
\qfigdef{p1_chem_s1_q03_B}{Score Advanced - Test 21 (08 Apr 2024)-img/p1_chem_s1_q03_B.png}{55.6}{0}
\qfigdef{p1_chem_s1_q03_C}{Score Advanced - Test 21 (08 Apr 2024)-img/p1_chem_s1_q03_C.png}{42.3}{0}
\qfigdef{p1_chem_s1_q03_D}{Score Advanced - Test 21 (08 Apr 2024)-img/p1_chem_s1_q03_D.png}{46.1}{0}
\qfigdef{p1_chem_s1_q04_1}{Score Advanced - Test 21 (08 Apr 2024)-img/p1_chem_s1_q04_1.png}{92.1}{0}
\qfigdef{p1_chem_s2_q09_1}{Score Advanced - Test 21 (08 Apr 2024)-img/p1_chem_s2_q09_1.png}{139.6}{0}
\qfigdef{p1_chem_s2_q09_A}{Score Advanced - Test 21 (08 Apr 2024)-img/p1_chem_s2_q09_A.png}{106.3}{0}
\qfigdef{p1_chem_s2_q09_B}{Score Advanced - Test 21 (08 Apr 2024)-img/p1_chem_s2_q09_B.png}{107.2}{0}
\qfigdef{p1_chem_s2_q09_C}{Score Advanced - Test 21 (08 Apr 2024)-img/p1_chem_s2_q09_C.png}{71.0}{0}
\qfigdef{p1_chem_s2_q09_D}{Score Advanced - Test 21 (08 Apr 2024)-img/p1_chem_s2_q09_D.png}{82.1}{0}
\qfigdef{p1_chem_s2_q10_1}{Score Advanced - Test 21 (08 Apr 2024)-img/p1_chem_s2_q10_1.png}{112.5}{0}
\qfigdef{p1_chem_s2_q10_A}{Score Advanced - Test 21 (08 Apr 2024)-img/p1_chem_s2_q10_A.png}{65.1}{0}
\qfigdef{p1_chem_s2_q10_B}{Score Advanced - Test 21 (08 Apr 2024)-img/p1_chem_s2_q10_B.png}{65.8}{0}
\qfigdef{p1_chem_s2_q10_C}{Score Advanced - Test 21 (08 Apr 2024)-img/p1_chem_s2_q10_C.png}{67.6}{0}
\qfigdef{p1_chem_s2_q10_D}{Score Advanced - Test 21 (08 Apr 2024)-img/p1_chem_s2_q10_D.png}{59.9}{0}
\qfigdef{p1_chem_s3_para5_6_1}{Score Advanced - Test 21 (08 Apr 2024)-img/p1_chem_s3_para5_6_1.png}{245.2}{0}
\raggedcolumns
\makeatletter\@beginparpenalty=10000\makeatother
\newcolumntype{L}[1]{>{\raggedright\arraybackslash}p{#1}}
\begin{document}
\begin{center}{\Large\bfseries Score Advanced --- Test 21 (08 Apr 2024)}\end{center}
\begin{multicols}{2}
\raggedcolumns
\papersec{1}
\partsec{1}{PHYSICS}
\sect{SECTION-I (i)}
\qn{1} Three capacitor $C_1=6$F, $C_2=3$F and $C_3=4$F are separately charged with battery such that charge on each capacitor are 36C, 24C and 40C respectively and then disconnected with the battery. After that they are connected in a circuit along with an uncharged capacitor $C_4=4$F and a battery of emf $=4$ volts as shown in figure. Now the switch S is closed. Choose the correct option(s).\par
\qfig{p1_phy_s1_q01_1}
\begin{qoptions}
\opt{A}{Final charge on capacitor $C_4=4$C}
\opt{B}{Final charge on capacitor $C_2=24$C}
\opt{C}{Final charge on capacitor $C_3=36$C}
\opt{D}{Amount of charge flow in circuit $=$ zero}
\end{qoptions}
\qn{2} Two hemispherical thin shells made of insulating materials are concentrically arranged in a free space as shown. The radii of the smaller and larger hemispheres are 'R' and '2R' and they carry positive charges 'Q' and '2Q' respectively. The charges are uniformly distributed over the surfaces of the shells. Then choose the correct option(s)\par
\qfig{p1_phy_s1_q02_1}
\begin{qoptions}
\opt{A}{The electrostatic force between the hemispherical shell is $\dfrac{Q^2}{8\pi\varepsilon_oR^2}$}
\opt{B}{The electrostatic force between the hemispherical shell is $\dfrac{Q^2}{16\pi\varepsilon_oR^2}$}
\opt{C}{The net electric potential at the centre 'O' of the shells is $\dfrac{Q}{2\pi\varepsilon_oR}$}
\opt{D}{The net electric potential at the centre 'O' of the shells is $\dfrac{3Q}{2\pi\varepsilon_oR}$}
\end{qoptions}
\qn{3} A uniform solid ball of radius $R$ rolls without slipping between the rails such that the horizontal distance is $\sqrt3R$ between the two contact points of the rails with the ball. Figure (a) shows the front view of the ball and figure (b) shows the side view of the ball. $V_{CM}$ is the velocity of centre of mass of the ball and $\omega$ is the angular velocity of the ball after rolling down a distance $2h$ along the incline then\par
\qfig{p1_phy_s1_q03_1}
\begin{qoptions}
\opt{A}{$V_{CM}=\omega R$}
\opt{B}{$V_{CM}=\omega\dfrac{R}{2}$}
\opt{C}{$V_{CM}=\sqrt{\dfrac{10gh}{13}}$}
\opt{D}{$V_{CM}=\sqrt{\dfrac{10gh}{7}}$}
\end{qoptions}
\qn{4} $n$ number of identical equilateral prism are kept in contact as shown in figure. If deviation through a single prism is $\delta$. Then ($n$, $m$ are integers)(Assume surrounding medium be air)\par
\qfig{p1_phy_s1_q04_1}
\begin{qoptions}
\opt{A}{If $n=2m$, deviation through $n$ prism is zero.}
\opt{B}{If $n=2m+1$, deviation through system of $n$ prisms is $\delta$.}
\opt{C}{If $n=2m$, deviation through system of $n$ prisms is $\delta$.}
\opt{D}{If $n=2m+1$, deviation through system of $n$ prism is zero.}
\end{qoptions}
\qn{5} A partition divides a container having insulated walls into two compartments I and II. The same gas fills the two compartments whose initial parameters are given. The partition is a conducting wall which can move freely without friction. Which of the following statement's is/are correct, with reference to the final equilibrium position?\par
(initially the partition is held at rest)\par
\qfig{p1_phy_s1_q05_1}
\begin{qoptions}
\opt{A}{The Pressure in the two compartments are equal.}
\opt{B}{Volume of compartment I is $\dfrac{3V}{5}$}
\opt{C}{Volume of compartment II is $\dfrac{12V}{5}$}
\opt{D}{Final pressure in compartment I is $\dfrac{5P}{3}$}
\end{qoptions}
\qn{6} A double slit setup is shown in the figure. One of the slits is in medium 2 of refractive index $n_2$. The other slit is at the interface of this medium with another medium 1 of refractive index $n_1$ ($\neq n_2$). The line joining the slits is perpendicular to the interface and the distance between the slits is $d$. The slit widths are much smaller than $d$. A monochromatic parallel beam of light is incident on the slits from medium 1. A detector is placed in medium 2 at a large distance from the slits, and at an angle $\theta$ from the line joining them, so that $\theta$ equals the angle of refraction of the beam. Consider two approximately parallel rays from the slits received by the detector.\par
\qfig[0.85]{p1_phy_s1_q06_1}
Which of the following statement(s) is/are correct ?
\begin{qoptions}
\opt{A}{The phase difference between the two rays is independent of $d$.}
\opt{B}{The two rays interfere constructively at the detector.}
\opt{C}{The phase difference between the two rays depends on $n_1$ but is independent of $n_2$.}
\opt{D}{The phase difference between the two rays vanishes only for certain values of $d$ and the angle of incidence of the beam, with $\theta$ being the corresponding angle of refraction.}
\end{qoptions}
\sect{SECTION-I (ii)}
\paraq{7}{8}A non conducting semi circular ring of radius $R$ having charge distributed over its circumference according to the formula $\lambda=\lambda_0\sin\theta$ is placed in a uniform electric field $E$ in such a way that plane of the semi circular ring is parallel to electric field $E$ as shown in figure. Consider that polarizability of the semi circular ring is equal to zero.\par
\qfig{p1_phy_s2_para7_8_1}
\qn{7} Potential at the midpoint of line AB in the absence of electric field $E$ (consider that potential at infinite is zero).
\begin{qoptions}
\opt{A}{$\dfrac{\lambda_o}{\pi\varepsilon_o}$}
\opt{B}{$\dfrac{\lambda_o}{2\pi\varepsilon_o}$}
\opt{C}{$\dfrac{2\lambda_o}{\pi\varepsilon_o}$}
\opt{D}{$\dfrac{\lambda_o}{4\pi\varepsilon_o}$}
\end{qoptions}
\qn{8} Potential difference between point A and point B in presence of electric field $E$
\begin{qoptions}
\opt{A}{$\dfrac{\lambda_o}{2\pi\varepsilon_o}+E.R$}
\opt{B}{$\dfrac{\lambda}{\pi\varepsilon_o}+E.R$}
\opt{C}{$\dfrac{\lambda}{2\pi\varepsilon_o}$}
\opt{D}{$E.R$}
\end{qoptions}
\paraq{9}{10}A special metal S conducts electricity without any resistance. A closed wire loop, made of S, does not allow any change in flux through itself by inducing a suitable current to generate a compensating flux. The induced current in the loop cannot decay due to its zero resistance. This current give rise to a magnetic moment which in turn repels the source of magnetic field or flux. Consider such a loop, of radius $a$, with its center at the origin. A magnetic dipole of moment $m$ is brought along the axis of this loop from infinity to a point at distance $r$ ($\gg a$) from the center of the loop with its north pole always facing the loop, as shown in the figure below.\par
The magnitude of magnetic field of a dipole $m$, at a point on its axis at distance $r$, is $\dfrac{\mu_om}{2\pi r^3}$, where $\mu_o$, is the permeability of free space. The magnitude of the force between two magnetic dipoles with moments, $m_1$ and $m_2$, separated by a distance $r$ on the common axis, with their north poles facing each other, is $\dfrac{km_1m_2}{r^4}$, where $k$ is a constant of appropriate dimensions. The direction of this force is along the line joining the two dipoles.\par
\qfig{p1_phy_s2_para9_10_1}
\qn{9} When the dipole $m$ is placed at a distance $r$ from the center of the loop (as shown in the figure), the current induced in the loop will be proportional to
\begin{qoptions}
\opt{A}{$\dfrac{m}{r^3}$}
\opt{B}{$\dfrac{m^2}{r^2}$}
\opt{C}{$\dfrac{m}{r^2}$}
\opt{D}{$\dfrac{m^2}{r}$}
\end{qoptions}
\qn{10} The work done in bringing the dipole slowly from infinity to a distance $r$ from the center of the loop by the given process is proportional to
\begin{qoptions}
\opt{A}{$\dfrac{m}{r^5}$}
\opt{B}{$\dfrac{m^2}{r^5}$}
\opt{C}{$\dfrac{m^2}{r^6}$}
\opt{D}{$\dfrac{m^2}{r^7}$}
\end{qoptions}
\sect{SECTION-II (i)}
\paraq{1}{2}A sleeve of mass $m_1=5$ kg can slide on a smooth horizontal rod. A block of mass $m_2=2$ kg is suspended at the mid-point of a light inextensible string of length $\ell=2.3$ m, one end of which is attached to the sleeve and the other end is attached to one end of the rod. Initially the sleeve is held at rest and the angle between the rod and a string segment is $\theta_0=53^\circ$ as shown in the figure. A constant horizontal force $F=20$N is applied on the sleeve. When the angle between the rod and the string segment becomes $\theta=37^\circ$, the speeds of the sleeve and the block are $v_1$ m/s and $v_2$ m/s respectively. (Take $g=10$ m/s$^2$)\par
\qfig{p1_phy_s3_para1_2_1}
\qn{1} The value of $v_1$ is \underline{\hspace{1.5cm}}
\qn{2} The value of $v_2$ is \underline{\hspace{1.5cm}}
\paraq{3}{4}In a stationary wagon of mass $M=32$ kg and length $d=1$m. a small block of mass $m=1$ kg is projected along its length with a velocity $v=10$ m/s towards right end from left end as shown in the figure. If coefficient of restitution for all collision is $e=\dfrac12$. All contact surfaces are smooth and each collision is instantaneous.\par
\qfig{p1_phy_s3_para3_4_1}
\qn{3} Velocity (in m/s) of the wagon after two collisions is \underline{\hspace{1.5cm}}
\qn{4} Time (in sec) after the projection of the block when the block comes at rest with respect to ground is \underline{\hspace{1.5cm}}
\paraq{5}{6}A block of ice at $0^\circ$C and mass $M$ kg is kept inside isolated container of negligible heat capacity. Steam of mass $m$ kg at temperature $100^\circ$C is kept inside the container. The ice melted and its temperature becomes $20^\circ$C in steady state. The steam of mass 0.2 kg at $100^\circ$C is again let inside the container, due to which temperature inside the container becomes $30^\circ$C in steady state. Latent heat of ice and steam are 80 kcal/kg and 540 kcal/kg respectively. Specific heat capacity of water is $S_w=1$ kcal/kg$^\circ$C.
\qn{5} The value of $m$ (in kg) is \underline{\hspace{1.5cm}}
\qn{6} The value of $M$ (in Kg) is \underline{\hspace{1.5cm}}
\sect{SECTION-II (ii)}
\qn{7} Consider a hydrogen-like ionized atom with atomic number $Z$ with a single electron. In the emission spectrum of this atom the photon emitted in the $n=2$ to $n=1$ transition has energy 74.8 eV higher than the photon emitted in the $n=3$ to $n=2$ transition. The ionization energy of the hydrogen atom is 13.6 eV. The value of $Z$ is\underline{\hspace{1cm}}
\qn{8} The figure shows a portion of an electric circuit. Resistors are known and are indicated on the diagram and the voltmeters are identical. If the voltmeters $V_1$ and $V_2$ read 7.5V and 5.0V respectively, find the reading of voltmeter $V_3$ (in volts).\par
\qfig{p1_phy_s4_q08_1}
\qn{9} A current $I=20$ ampere flows through a wire shaped in the form of an infinite parabola of latus rectum $4a$ ($a=1$ mm). If the magnetic field at the focus of the parabola is $\eta\pi\times10^{-\beta}$ tesla, where $\eta$ and $\beta$ are the positive integers, less than 9. Find the value of $\eta\times\beta$.\par
\qfig{p1_phy_s4_q09_1}
\partsec{2}{CHEMISTRY}
\sect{SECTION-I (i)}
\qn{1} Which of the following statements is/are not correct for following compounds\par
(I) SCl$_2$ (OCH$_3$)$_2$ and (II) SF$_2$ (OCH$_3$)$_2$
\begin{qoptions}
\opt{A}{OCH$_3$ groups in both cases occupy the same position}
\opt{B}{Cl $-$ atoms occupy equatorial position in case of (I) and F$-$ atoms occupy equatorial position in case of (II)}
\opt{C}{Cl $-$ atoms occupy axial position in case of (I) and F $-$ atoms occupy equatorial position in case of (II)}
\opt{D}{Cl and F atoms all occupy axial position in case of both (I) and (II)}
\end{qoptions}
\qn{2} Identify the incorrect statement(s) :-
\begin{qoptions}
\opt{A}{Combustion of methane in an adiabatic flask will cause no change in temperature of the system}
\opt{B}{Dimerization of acetic acid causes an increase in entropy of the system}
\opt{C}{It is possible to have both adiabatic reversible and adiabatic irreversible processes between two states}
\opt{D}{P - V work is always non-zero when there is a change in volume}
\end{qoptions}
\qn{3} Which compounds would not give 5-keto-2-methylhexanal upon ozonolysis?
\begin{qoptions}
\opt{A}{\qfigi{p1_chem_s1_q03_A}}
\opt{B}{\qfigi{p1_chem_s1_q03_B}}
\opt{C}{\qfigi{p1_chem_s1_q03_C}}
\opt{D}{\qfigi{p1_chem_s1_q03_D}}
\end{qoptions}
\qn{4} Correct statement(s) relative to given base is/are\par
\qfig{p1_chem_s1_q04_1}
\begin{qoptions}
\opt{A}{Basic strength at (T) is greater than (S)}
\opt{B}{pK$_b$ difference between (R) and (T) is greater (in magnitude) than the difference between (P) and (T)}
\opt{C}{Hydrogen from nitrogen at (P) is more acidic than hydrogen from nitrogen at (Q)}
\opt{D}{Delocalization of lone pair on nitrogen at (S) will bring aromaticity to the smaller ring}
\end{qoptions}
\qn{5} If the radius of first Bohr's orbit of H-atom is $x$, which of the following is/are the correct conclusion(s)?
\begin{qoptions}
\opt{A}{The de-Broglie wavelength in the third Bohr orbit of H-atom$=6\pi x$}
\opt{B}{The fourth Bohr's radius of He$^+$ ion $=8x$}
\opt{C}{The de-Broglie wavelength in third Bohr's orbit of Li$^{2+}=2\pi x$}
\opt{D}{The second Bohr's radius of Be$^{3+}=x$}
\end{qoptions}
\qn{6} For the following options, select all the right combinations of concentration cell and condition for the spontaneity given $P_1$ and $P_2$ are gas pressures and $C_1$, $C_2$ are concentration terms with all other tern carrying their usual meanings
\begin{qoptions}
\opt{A}{$\mathrm{Cl^-(C_1)/\underset{1\,atm}{Cl_2},Pt\,||\,Pt,\underset{1\,atm}{Cl_2}/Cl^-(C_2)};\ C_1>C_2$}
\opt{B}{$\mathrm{Fe_{(s)}\,|Fe^{+3}_{(aq)}(C_1)|\,|Fe^{+3}_{(aq)}(C_2)|\,Fe};\ C_1>C_2$}
\opt{C}{$\mathrm{Pt\,(\underset{P_1}{H_2})|HCl|Pt\,(\underset{P_2}{H_2})};\ P_1>P_2$}
\opt{D}{$\mathrm{Pd\,(\underset{1\,atm}{H_2})|HCl(C_1)|\,|HCl(C_2)|\,Pd\,(\underset{1\,atm}{H_2})};\ C_2>C_1$}
\end{qoptions}
\sect{SECTION-I (ii)}
\paraq{7}{8}A sample (X) of compound(s) contains any one or two of the following AgNO$_3$, ZnSO$_4$,CuSO$_4$, MgCl$_2$, CaCO$_3$, BaCl$_2$, Na$_2$SO$_4$, NaCl, NaOH\par
When (X) is dissolved in water, a clear solution is obtained. This solution is divided into two parts I and II. Part I turns pink on addition of a few drops of phenolphthalein. When dilute hydrochloric acid is added slowly to part II, a white precipitate (Y) is obtained which dissolves on further addition of dilute hydrochloric acid
\qn{7} The clear solution obtained when (X) is dissolved in water, contains
\begin{qoptions}
\opt{A}{ZnCl$_4^{2-}$}
\opt{B}{AgCl$_2^{-}$}
\opt{C}{Zn(OH)$_4^{2-}$}
\opt{D}{CuCl$_4^{2-}$}
\end{qoptions}
\qn{8} The sample (X) contains
\begin{qoptions}
\opt{A}{AgNO$_3$ \& NaCl}
\opt{B}{ZnSO$_4$ \& BaCl$_2$}
\opt{C}{HgCl$_2$ \& NaOH}
\opt{D}{ZnSO$_4$ \& NaOH}
\end{qoptions}
\paraq{9}{10}Due to severe angle strain, epoxides undergo ring opening reactions readily. In the presence of an acid catalyst, such as sulfuric acid, epoxides are hydrolyzed to glycols. Because there is some carbocation character developed in the transition state for an acid-catalyzed epoxide ring opening, attack of the nucleophile on unsymmetrical epoxides occurs preferentially at the carbon better able to bear a partial positive charge. Epoxides undergo ring-opening reactions with a variety of nucleophiles. Good nucleophiles attack an epoxide ring by an S$_N$2 mechanism and show an S$_N$2-like regioselectivity; that is, the nucleophile attacks at the less hindered carbon.
\qn{9} The product of the following reaction is\par
\qfig{p1_chem_s2_q09_1}
\begin{qoptions}
\opt{A}{\qfigi{p1_chem_s2_q09_A}}
\opt{B}{\qfigi{p1_chem_s2_q09_B}}
\opt{C}{\qfigi{p1_chem_s2_q09_C}}
\opt{D}{\qfigi{p1_chem_s2_q09_D}}
\end{qoptions}
\qn{10} The product of the following reaction is\par
\qfig{p1_chem_s2_q10_1}
\begin{qoptions}
\opt{A}{\qfigi{p1_chem_s2_q10_A}}
\opt{B}{\qfigi{p1_chem_s2_q10_B}}
\opt{C}{\qfigi{p1_chem_s2_q10_C}}
\opt{D}{\qfigi{p1_chem_s2_q10_D}}
\end{qoptions}
\sect{SECTION-II (i)}
\paraq{1}{2}An enclosed glass bulb of volume '$V$' contains a pinch of solid NH$_4$Cl and $\dfrac{1}{100}$ moles of Ne gas. At 300 K, the internal pressure inside the glass bulb is 114mm Hg which increases to 908 mm Hg when the bulb is heated to 600 K. (Assume ideal gas behavior, no$-$sublimation and complete dissociation of NH$_4$Cl at 600K)\par
[ The only reaction taking place is : NH$_4$Cl$_{(s)}\rightarrow$ NH$_{3(g)}$ + HCl$_{(g)}$ ]
\qn{1} What will be the partial pressure of NH$_3$(g) (in mm of Hg) in the bulb at 600K ?
\qn{2} What mass of NH$_4$Cl (in grams) is approximately present in 'Pinch' ?
\paraq{3}{4}NaClO$_3$ on reaction with SO$_2$ gas is presence of H$_2$SO$_4$ and trace of NaCl forms a yellow gas (A) as one of the product. Ozonolysis of 'A' produces 'X' along with liberation of oxygen gas.
\qn{3} The oxidation state of central atom in 'A' is '$+y$'. value of '$y$' is.
\qn{4} 1 mole of 'X' on reaction with excess of NaOH gives 'B', 'Q' and H$_2$O. Oxidation states of central atoms in the anions of 'B' and 'Q' and '$+p$' and '$+q$' respectively. Value of $|q-p|$ is
\paraq{5}{6}Phenol is one of the starting material of two commercially important compounds "X" \& "Y". They are prepared as follows.\par
\qfig{p1_chem_s3_para5_6_1}
\qn{5} What is the mass (in grams) of 'X' obtained in Experiment - I
\qn{6} What is the mass (in grams) of "Y" obtained in Experiment - II
\sect{SECTION-II (ii)}
\qn{7} In a vessel, two equilibria are simultaneously established at a given temperature as follows:\par
N$_2$(g) + 3H$_2$(g) $\rightleftharpoons$ 2NH$_3$ (g) $[K_{P_1}]$\par
N$_2$(g) + 2H$_2$(g) $\rightleftharpoons$ N$_2$H$_4$ (g) $[K_{P_2}]$\par
Initially the vessel contains N$_2$ and H$_2$ in the molar ratio of 9 : 13. The equilibrium pressure is $7P_0$. The pressure due to NH$_3$ is $P_0$ and due to H$_2$ is $2P_0$. If the value of $(K_{P_1}+K_{P_2})$ is $\dfrac{1}{aP_0^2}$.\par
Find the value of $a$
\qn{8} If the freezing point of 0.01 molal aqueous solution of a cobalt (III) chloride - ammonia complex (which behaves as a strong electrolyte) is $-0.0558^\circ$C, the number of chloride(s) in the coordination sphere of the complex is [$K_f$ of water $=1.86$ K kg mol$^{-1}$]
\qn{9} The number of chloride ions surrounding each Fe$^{+3}$ ion in solid anhydrous ferric chloride is
\partsec{3}{MATHEMATICS}
\sect{SECTION-I (i)}
\qn{1} Let $\displaystyle\int e^{x^2}.e^x\left(2x^2+x+1\right)dx=e^{x^2}.f(x)+C$ where $f(x)$ is some non-zero constant function and $C$ is some arbitrary constant. If the local minimum value of $f(x)$ is equal to $m$, then:\par
[[.] denotes greatest integer function]
\begin{qoptions}
\opt{A}{$f(x)$ is increasing in $(0,\infty)$}
\opt{B}{the value of $\displaystyle\lim_{x\to0}(1+f(x))^{1/x}$ is equal to 1}
\opt{C}{the value of $\displaystyle\int_0^1\left(f(x)+e^x\right)dx$ is equal to $2e$}
\opt{D}{the value of $\left[\dfrac{-1}{m}\right]$ is equal to 2.}
\end{qoptions}
\qn{2} Let $f_1(x)=\sin\pi x$, $f_2(x)=\cos\pi x$, $f_3(x)=[x]$, $f_4(x)=\{x\}$, $f_5(x)=-2x$, $f_6(x)=|x|$ be given functions then which of the following statements is/are not correct? (Where [.] is G.I.F and \{.\} is F.P.F., $|x|$)
\begin{qoptions}
\opt{A}{Two functions $f_a(x)$ and $f_b(x)$ are selected at random then the probability that $f_a(f_b(x))$ is a constant function ($a\neq b$) is ($a=1,2,\ldots..6$ ; $b=1,2,\ldots.6$) $\dfrac15$.}
\opt{B}{Two functions are selected at random and if it is known that their sum is a periodic, function, probability that the period is 2 is $\dfrac15$.}
\opt{C}{If a function is selected at random from the given functions then the probability that it is a bijection from R to R is $\dfrac34$}
\opt{D}{If two functions are selected at random and added then the probability that the sum is an even function is $\dfrac{1}{15}$}
\end{qoptions}
\qn{3} The function $f$ is continuous and has the property $f(f(x))=1-x$ for all $x\in R$ and $J=\displaystyle\int_0^1f(x)\,dx$ then
\begin{qoptions}
\opt{A}{$f\left(\dfrac14\right)+f\left(\dfrac34\right)=1$}
\opt{B}{The value of $J$ equals to $\dfrac12$}
\opt{C}{$f\left(\dfrac13\right)f\left(\dfrac23\right)=1$}
\opt{D}{$\displaystyle\int_0^{\pi/2}\dfrac{\sin x}{(\sin x+\cos x)^3}dx$ has the same value as $J$}
\end{qoptions}
\qn{4} If $a,b,c\in R^+$ then $\displaystyle\lim_{n\to\infty}\sum_{k=1}^n\dfrac{n}{(an+k)(bn+k)}=$
\begin{qoptions}
\opt{A}{$\dfrac{1}{a-b}\log\dfrac{b(a+1)}{a(a+1)}$ if $a\neq b$}
\opt{B}{$\dfrac{1}{a(a+1)}$ if $a=b$}
\opt{C}{$\dfrac{1}{a-b}\log\dfrac{a(b+1)}{b(a+1)}$ if $a\neq b$}
\opt{D}{does not exists if $a=b$}
\end{qoptions}
\qn{5} Consider the two parabolas $y^2=x-1$ and $x=y-1$. Let $A$ be the area between the coordinate axes, the two parabolas and the line which makes the shortest distance between two parabolas and $A^1$ be the area between $x=0$, $x^2=y-1$, $x=y$ and the line of shortest distance between $y^2=x-1$ and $x^2=y-1$ then which of the following is/are true?
\begin{qoptions}
\opt{A}{$A=A^1$}
\opt{B}{$A=2A^1$}
\opt{C}{The shortest distance between the parabolas is $\dfrac{3}{2\sqrt2}$}
\opt{D}{The shortest distance between the parabolas is $\dfrac12$}
\end{qoptions}
\qn{6} Let $f(x)=x^3-3x+1$ then which of the following can be a solution of the equation $f(x)=0$
\begin{qoptions}
\opt{A}{$2\cos\dfrac{2\pi}{9}$}
\opt{B}{$2\cos\dfrac{4\pi}{9}$}
\opt{C}{$2\cos\dfrac{8\pi}{9}$}
\opt{D}{$2\sin\dfrac{\pi}{18}$}
\end{qoptions}
\sect{SECTION-I (ii)}
\paraq{7}{8}Let $a_1,a_2,a_3\ldots..$be an arithmetic sequence and $b_1,b_2,b_3\ldots..$ be a geometric sequence. Suppose $a_1b_1=20$, $a_2b_2=19$, $a_3b_3=14$. Now the value of $a_4b_4$ needs to be maximized.
\qn{7} The common ratio of the geometric sequence to achieve the maximization is
\begin{qoptions}
\opt{A}{$\dfrac32$}
\opt{B}{$\dfrac12$}
\opt{C}{$\dfrac75$}
\opt{D}{$\dfrac13$}
\end{qoptions}
\qn{8} The maximum value of $a_4b_4$ is
\begin{qoptions}
\opt{A}{1.25}
\opt{B}{8.50}
\opt{C}{9.25}
\opt{D}{9.75}
\end{qoptions}
\paraq{9}{10}A ray of light emanating from a point A($1,2,3$) and strikes a line minor $L:\dfrac{x+2}{1}=\dfrac{y-3}{2}=\dfrac{z-1}{3}$ at point P, after reflection it passes through the point B($3,6,9$).
\qn{9} The equation of plane passing through the points P, A and B is
\begin{qoptions}
\opt{A}{$x+y-z=0$}
\opt{B}{$2x-y+z=9$}
\opt{C}{$5x-7y+3z=0$}
\opt{D}{$4x-5y+2z=0$}
\end{qoptions}
\qn{10} The value of $PA+PB$ equals
\begin{qoptions}
\opt{A}{$\dfrac{7}{\sqrt2}$}
\opt{B}{$7\sqrt2$}
\opt{C}{$\dfrac{49}{\sqrt2}$}
\opt{D}{$49\sqrt2$}
\end{qoptions}
\sect{SECTION-II (i)}
\paraq{1}{2}$A_n$ and $B_n$ be square matrices of order 3 which are defined as $A_n=[a_{ij}]$ and $B_n=[b_{ij}]$ where $a_{ij}=\dfrac{2i+j}{3^{2n}}$ and $b_{ij}=\dfrac{3i-j}{2^{2n}}$ $\forall\,i,j,1\le i,\ j\le3$. If $L=\displaystyle\lim_{n\to\infty}\mathrm{Tr}\left(3A_1+3^2A_2+3^3.A_3+\ldots..+3^n.A_n\right)$,\par
$M=\displaystyle\lim_{n\to\infty}\mathrm{Tr}\left(2B_1+2^2B_2+2^3.B_3+\ldots..+2^n.B_n\right)$\par
(Where Tr (P) = Trace of matrix P). Let $p,n\in N$,\par
$f(x)=1-x^p$ and\par $g_n(x)=\dfrac{n}{\dfrac{1}{f(x)}+\dfrac{1}{f(2x)}+\ldots\ldots+\dfrac{1}{f(nx)}}$ and\par $R=\displaystyle\lim_{x\to0}\dfrac{1-g_n(x)}{x^p}$\par
at $n=5$ and $p=3$
\qn{1} The value of $100-R$ is equal to \underline{\hspace{1.5cm}}
\qn{2} The value of $L+M+R$ is equal to
\paraq{3}{4}Let $f_1(x)=|x-1|+|x-2|+|x-3|$,\par $f_2(x)=|x-1|+|x-2|+|x-3|+|x-4|$, $A$ = minimum value of $f_1(x)$,\par
$B$ = minimum value of $f_2(x)$, $f(x)=12x^2\displaystyle\int_0^{A/2}yf(y)dy+20\int_0^{B/4}xy^2f(y)dy+4x$
\qn{3} The maximum value of $f(x)$, is for $x\in$ Real\underline{\hspace{1.5cm}}
\qn{4} Area bounded by the curve $y=f(x)$ and x - axis is \underline{\hspace{1.5cm}}
\paraq{5}{6}If $p\neq1$ is an $n^{th}$ root of unity having least argument then\par
$1+5p+9p^2+\ldots+(4n-3)p^{n-1}=\dfrac{an}{p-1}$ and hence the values of the sums\par
$1+5\cos\left(\dfrac{2\pi}{n}\right)+9\cos\left(\dfrac{4\pi}{n}\right)+\ldots+(4n-3)\cos\left(\dfrac{2(n-1)\pi}{n}\right)=b.n$ and\par
\resizebox{\linewidth}{!}{$5\sin\left(\dfrac{2\pi}{n}\right)+9\sin\left(\dfrac{4\pi}{n}\right)+\ldots+(4n-3)\sin\left(\dfrac{2(n-1)\pi}{n}\right)=c\left(n\cot\dfrac{\pi}{n}\right)$}\par for $n\ge2$
\qn{5} The value of $a+b+c$ is equal to \underline{\hspace{1.5cm}}
\qn{6} The argument of $z=a-i(b+c)$ is $\theta$ ($\theta$ is in radians) then $8\theta=$ \underline{\hspace{1.5cm}}
\sect{SECTION-II (ii)}
\qn{7} Let $g(\theta)=\displaystyle\int_{-(\hat a\cdot\hat b)^2}^{|\hat a\times\hat b|^2}(2t+1)\,dt$ where $\theta$ is the angle between $\hat a$ and $\hat b$. If volume of the parallelopiped whose coterminous edges are represented by the vectors $\hat a$, $\hat a\times\hat b$ and $\hat a\times\left(\hat a\times\hat b\right)$ (Where angle between $\hat a$ and $\hat b$ is taken from the equation $2g(\theta)-1=0$) is $\dfrac pq$ then the least value of $p+q$ is \underline{\hspace{1cm}} mention $\hat a$ \& $\hat b$ are unit vectors where $p,q\in N$
\qn{8} Let $P_k=\begin{cases}k,&\text{if }k\text{ is odd}\\[2pt]\dfrac{1}{k-1},&\text{if }k\text{ is even}\end{cases}$ and $\displaystyle\prod_{k=1}^{14}\left(P_K\log_{(k+1)}(k+2)\right)=N$, then the number of solutions of the equation $\tan^2x+\cot^2x+2=N\sin^22x$ in $x\in[0,2\pi]$ is \underline{\hspace{1.5cm}}
\qn{9} Let $f(x)=x+x^2+x^4+x^8+x^{16}+x^{32}+\ldots.$ If the coefficient of $x^{10}$ in $f(f(x))$ is $k$, then $\dfrac k5=$ \underline{\hspace{1cm}}
\end{multicols}
\end{document}
